% !TEX root = ../thesis-example.tex
%
\chapter{Introduction}
\label{chap:introduction}

Brain Computer Interfaces (BCIs), first used in experiments in the year 1969, aim to record the electrical activity inside the brain and influence computers and/or machines with it, while also providing the possibility of having a direct communication path between the brain and its external environment \cite{kawala-sterniuk-2021}.
While BCIs are probably primarily known from clinical context like neural rehabilitation and assistive technologies like controlling prostheses, according to a review by Värbu et al. published in 2022, medical applications only amount for about 30\% inside the articles published about BCIs between the year of 2009 and 2019 \cite{varbu-2022}. Most commonly, BCIs are used for non-clinical tasks, primarily for monitoring tasks, such as emotion recognition \cite{yang2018continuous} or mental fatigue estimation \cite{tian2018mental}.
Another non-clinical application for BCIs can be found in the area of gaming, where BCIs can be used to control parts of video games, for example by tracking brain signal frequencies or actively detecting different kinds of spikes in brain activity. 
% maybe go more into the motor imaging stuff already here?
Another way of using BCIs in video games is by using so-called \grqq mental commands\grqq %cite!
, which are based on the process of motor imaging and where the user actively thinks about a self-chosen, mentally accurately reproducible task which the software of the BCI records. If a user then thinks about said self-chosen task, the BCI is able to detect that and is then able to act on this command by moving a video game character forward or make it use an item, to stay in the context of video games. 
However, in order for the BCI to accurately and efficiently detect those mental commands, the user needs to train the BCI on the mental command by thinking about the same task for the same command in a multitude of recording sessions. % maybe umformulieren
This process takes a lot of time and is further linearly increased with the amount of mental commands that the user wants to train the BCI to detect. In a context of using a BCI as part of the input for a video game, the user of such a video game would be required to undergo a lengthy training session of training mental commands prior to actually playing the video game in order to ensure an accurate detection of the mental commands in the video game. As the duration of the calibration process could directly and starkly negatively impact the player experience, the aim of this thesis is to analyse algortihms that aim to reduce this calibration time to enable players to start playing their game and use their BCI controller inputs more quicly. % possibly rephrase + add some quotations? Maaaaybe also slightly make the setion longer by talking about the motor imaging or briefly mention the BCI paradigms? 

% schauen, dass wir einheitlich entweder user oder player schreiben!



\section{Overview}

% can write this at a much later point in time

To understand the principles of Brain Computer Interfaces and how they work, relevant terms and definitions will be explained in chapter \ref{chap:definitions}, while giving direct examples where possible.
Then, in Chapter \ref{chap:applications}, three concrete examples for the use of BCIs in gaming application are presented, whilst trying to show a variety in the underlying systems that are used in them. chapter \ref{chap:challenges} shows and explain challenges that arise specifically in the use of BCI in gaming applications and suggests ways to overcome those. Finally, chapter \ref{chap:outlook} gives a possible outlook on future BCI applications in gaming before ultimately finishing the paper with a conclusion of the findings that were made throughout it.  